\documentclass[12pt]{article}
\usepackage[utf8]{inputenc}
\usepackage[T1]{fontenc}
\usepackage[margin=1in]{geometry}
\usepackage{setspace}
\usepackage{xcolor}
\usepackage[colorlinks=true, urlcolor=blue!60!black, linkcolor=blue!60!black]{hyperref}
\setlength{\parindent}{0em}
\setlength{\parskip}{0.6em}
\renewcommand{\baselinestretch}{1.0}
\begin{document}
\begin{center}
{\LARGE \textit{The Prize Cases}, 67 U.S. (2 Black) 635 (1863)\par}
\vspace{0.5em}
{\large Case Brief --- Two-Page Edition\par}
\end{center}

\section{Memory Jogger}
Upholding Lincoln's 1861 blockade, the Court held that the President cannot initiate war but must repel invasion or armed rebellion without waiting for Congress (\href{https://supreme.justia.com/cases/federal/us/67/635/}{67 U.S. 635}).

\section{Detailed Case Facts}
After Fort Sumter, with Congress in recess, Lincoln proclaimed a blockade of Southern ports in April 1861. Before Congress recognized the insurrection on July 13, 1861, Union ships captured four merchant vessels, Southern, British, and Mexican, as prizes. Congress retroactively approved Lincoln's acts on August 6, 1861.

\section{Procedural History}
Federal district courts condemned the vessels and cargoes; the claimants appealed directly, arguing that no legal war existed before July 13.

\section{Judicial Votes}
Affirmed, 5--4: Grier wrote for Wayne, Swayne, Miller, and Davis; Nelson dissented, joined by Taney, Catron, and Clifford.

\section{Holding}
The President, though he cannot declare war, is authorized and bound to resist invasion or rebellion by force, including blockade. Courts defer to his judgment that an insurrection has become war, and property of residents of hostile territory is capturable as enemies' property.

\section{Analysis of Opinions}
Grier reasoned that civil war is a fact, ``never solemnly declared,'' existing once rebels hold territory and close the courts; the President's executive and Commander-in-Chief powers and the 1795 and 1807 Acts required him to meet it. Even if legislative sanction was needed, the August 1861 Act cured any defect. Nelson distinguished war in a material sense from war in a legal sense, which only Congress can create, and held all captures before July 13 void; Lincoln could fight a ``personal war'' against insurgents but could not turn loyal citizens into enemies.

\section{Examples: Future Applications}
\textbf{Same-Side:} the President may repel a surprise attack on a naval base during a congressional recess. \textbf{Opposite-Side:} a preventive invasion of a nation that has not attacked exceeds his power. \textbf{Fence-Sitter:} retaliatory strikes after a crippling cyberattack on the power grid.

\section{Critique}
Originalists such as Michael Ramsey read the decision as confirming only a narrow defensive power, while John Yoo claims broad Article II war powers. Louis Fisher and others argue that modern presidents overread it as a blank check, ignoring that Lincoln acted defensively and Congress quickly ratified.

\section{Key Quotations}
\begin{quote}
``He has no power to initiate or declare a war either against a foreign nation or a domestic State. \ldots If a war be made by invasion of a foreign nation, the President is not only authorized but bound to resist force by force. He does not initiate the war, but is bound to accept the challenge without waiting for any special legislative authority.''
\end{quote}
\noindent --- Justice Grier, Opinion of the Court, 67 U.S. at 668.

\end{document}
